\section{Alat Desain dan Handoff ke Pengembang}

Untuk membuat wireframe dan mockup, tim sering menggunakan alat berbasis web atau desktop. \textbf{Figma} adalah salah satu alat populer yang memungkinkan kolaborasi daring, pembuatan komponen, dan prototype interaktif. Alternatif lain: Adobe XD, Sketch (macOS), Balsamiq, atau Lucidchart untuk diagram dan wireframe \cite{mdn-learn}.

Prinsip desain visual dasar yang perlu diperhatikan: \textbf{layout} (penataan elemen), \textbf{hierarki} (mana yang paling penting secara visual), dan \textbf{konsistensi} (warna, font, spasi yang seragam di seluruh halaman). Setelah mockup disetujui, \textbf{handoff} ke pengembang meliputi spesifikasi (ukuran, warna, font), aset (ikon, gambar), dan bila mungkin komponen yang dapat diinspeksi (misalnya kode CSS atau token desain) \cite{mdn-learn}.

Prototyping dan iterasi desain memungkinkan validasi dengan pengguna sebelum implementasi penuh, sehingga mengurangi biaya perbaikan di tahap pengembangan \cite{mdn-learn}.
